% Edition20, section12 opening: source collation and explicitly translator-owned proofs.
\egasection{§12开头：关联预概形、秩与两种投影的法英回校}

\noindent\textbf{实际范围与边界。}
本次直接查看法文手稿 \texttt{EGA5\_francais.pdf} 物理第64、65、66页及英文编辑稿 \texttt{EGAVaiello.pdf} 物理及印刷第43、44页的页面影像，并对第64页商模、转置句和第65页手添同构、Euler 式另作放大核对。有效正文从法文第64页“12. Généralisation de résultats antérieurs aux sections linéaires”开始，止于第66页第一大段末“pour $X^{(m)}$ sur $\Gr^m$”的右括号。紧接的“ La proposition 2.1. se transpose sans changement.”不属于本模块；它与英文第44页“The proposition 2.1 [se transpose Fr] translates (?) without change.”是下一模块的精确接点。

本模块对应旧 \texttt{v\_dot\_middle.tex} 的 DOT-MIDDLE BATCH 01--02。旧源全文保存在本版 \texttt{recollation/legacy\_v19/v\_dot\_middle.tex}；旧D17-M001至D17-M012的脚注原文又逐条保存于本报告末尾，便于核对原有解释与本次归因修订。独立第117批提供了候选译文、核验和例子，本次逐项回看源页并重作下面的局部论证，未把该批的旧核验记录视为本次已执行的检查。排版构建与最终输出检查由本版总检查记录另列，不由源页查看推定已经通过。

\subsection*{一、段落与源页映射}

\begin{longtable}{@{}p{.12\linewidth}p{.16\linewidth}p{.16\linewidth}p{.44\linewidth}@{}}
\toprule
正文标记 & 法文物理页 & 英文物理页 & 所保留内容\\
\midrule
\endhead
P01 & 64上 & 43中 & 标题、记号补充，$\operatorname{Grass}_n(P)$、$\operatorname{Grass}^n(P)$及空子簇约定。\\
P02 & 64上中 & 43中 & 先取秩 $r+1$ 的典范同构，随后重新取秩 $r$，余维 $m$ 参数。\\
P03 & 64中 & 43下 & 泛商、关联子函子、包含截面像的条件。\\
P04 & 64中下 & 43末至44首 & 任意 $T$、两个复合同态、转置与 $G_T(1)$ 的截面。\\
P05 & 65首 & 44首 & 零点预概形，记号 $H^{(m)}$，与第1节的 $m=1$ 情形相接。\\
P06 & 65上 & 44上 & $Y^{(m)}$、$Y^m_\xi$ 及原页手添的基底说明。\\
P07 & 65中 & 44中 & 作者的记号自评、迭代上标设想及试用说明。\\
P08 & 65末 & 44中 & 手添 $P$-同构、Grassmann 描述、核模与 Euler 式。\\
P09 & 66首 & 44中下 & 光滑、几何整及有理纤维；手添维数；原稿所说归约 $X=P$。\\
P10 & 66第一大段末 & 44下，2.1前 & 第1节图式类比，另一投影的射影丛及原括号内限制。\\
\bottomrule
\end{longtable}

标记P01--P10是源码中 \texttt{V20-12O-P01} 等注释的简写，并非另给作者的段落编号。各显示式的展开只显明原稿所说的复合与基变换；凡增加数学论证均放在译校注或本报告中。

\subsection*{二、错误归属与原稿改写层}

\begin{longtable}{@{}p{.18\linewidth}p{.33\linewidth}p{.40\linewidth}@{}}
\toprule
本版注号 & 源页实际证据 & 处理及与旧注的关系\\
\midrule
\endhead
D20-12O-01、02 & 法64的两处秩分别是 $r+1$ 和 $r$；固定 $m$ 处明确写余维。 & 不把这两处 $r$ 合并。英文的“given dimension $m$”是误译；D17-M002的归责限定于英文。\\
D20-12O-03 & 法64手添 $G$，后仍有“soit $F$”；底模排 $E_{\Gr}$，未见清楚的对偶号。 & 按已经明确的参数函子写 $E^\vee_{\Gr^m}\twoheadrightarrow G$。微小笔迹不被当作已恢复的对偶号；G/F冲突确在法文中。\\
D20-12O-04 & 法64两次乘积的底标为 $S$；截面与线性子簇均在 $S'$ 上。 & D17-M004所修的 $X$ 及 $S/S'$ 混排来自英文；法文没有相同错误。\\
D20-12O-05 & 法64文字目标为 $G_T^\vee$，但显示箭头及截面分别用 $G_T$、$G_T(1)$。 & 这一不协调来自法文，英译照存；正文订成 $G_T$，保留原错。\\
D20-12O-07 & 法65初次只写 $Y^m$，下句写 $Y^m_\xi$。 & 补齐下标并明示基变换；原稿手添“在 $P$ 上”“在 $P_{S'}$ 上”均保留。\\
D20-12O-08 & 法65迭代等式右端贴近裁边；英语给出完整等式。 & 保存作者的近似记号设想及可辨式，不宣布两个普遍族同构。旧D17-M008的旗标限制仍有效。\\
D20-12O-09 & 法65手添 $P$-；Euler 式手写小写 $f^*$。 & 恢复英译漏去的 $P$-，把旧D17-M009的 $F^*$ 混用归于英文；不再把法稿的 $f$ 判为大写 $F$。\\
D20-12O-10 & 法66手添 $m(r-m)$，英文相同。 & 这是法稿原错；继续采用旧译正确的 $m(r-1-m)$，并证明适用范围。\\
D20-12O-11、12 & 法66分别讨论 $H^{(m)}\to\Gr^m$ 与 $X^{(m)}\to\Gr^m$，并有明确限制括注。 & 保留两种投影的区别，不把对 $X$ 的光滑性转移到参数空间。\\
\bottomrule
\end{longtable}

法文第64页在“局部自由”后划去的短语为“de rang”（“秩为”）；在“soit $F$”之后、关联预概形句之前还划去了“Exploitation”（“应用／利用”）。这两个被删短语是稿件改写层，未重新当作有效正文。第65页 $F_X$ 后有覆打删字，第64、65页还可见零星覆写和校正符；其中不能可靠辨认的旧字没有补造。第65页可见的 $P$-插入、$f^*$、$F^\vee$、$\Omega^1_{P/S}$ 及第66页手添维数则已逐项核对。第65页靠近 Grassmann 显示式下方的数字1属于紧随其后的 $\Omega^1$ 上标，不能另读成一个第二的 Grassmann 下标。

英文第43页在秩 $r+1$ 后所加的“illegible, ask AG”作为英译编辑说明保留。法文的秩和维数数字均可辨，右缘连接缩写不被强行恢复为某个完整原词。第65页的迭代左式为 $X^{(m)(m')}$；右式可读至 $X^{(m+m}$，最后的撇号和括号受右裁边影响。英文相应写作
\[
X^{(m)(m_1)}=X^{m+m_1}.
\]
这只是作者自称“大体说来”的记号设想；本版没有把英文完整式反投射为法文每个被裁字符都已可辨。作者随后说“我们从未采用过这类约定”，法文使用过去的说法；本版据此译出，未沿用英文将其改成“以后绝不使用”的语气。

\subsection*{三、关联零点的概形结构：任意检验预概形上的核验}

本节以下各段均为译者核验，适用条件是 $S$ 为任意基预概形、$E$ 为有限秩 $r$ 的局部自由模。记
\[
\begin{gathered}
B=\Gr^m(P)=\operatorname{Grass}_m(E^\vee),\qquad E_B^\vee\twoheadrightarrow G,\\
N=G^\vee\hookrightarrow E_B,\qquad Q=E_B/N.
\end{gathered}
\]
在 $0\leq m\leq r$ 时，泛商列局部分裂，故 $N$ 与 $Q$ 局部自由，秩分别为 $m$ 与 $r-m$。这些事实由商的定义及局部自由商的局部分裂得到，不须假设基预概形约化或诺特。

对任意 $T\to P\times_S B$，以 $L=\mathcal O_T(1)$ 记 $P$ 因子给出的可逆商，并仍以 $N_T$ 记另一因子给出的子模。考虑
\[
\alpha:N_T\longrightarrow E_T\longrightarrow L.
\]
它的零点预概形由理想层
\[
\mathcal I_\alpha=
\Im\bigl(N_T\otimes L^\vee\longrightarrow\mathcal O_T\bigr)
\]
定义。选取局部基后，$\alpha$ 是一个有 $m$ 个系数的行向量；其转置 $L^\vee\to N_T^\vee=G_T$ 是同一组系数的列向量，而对应的 $G_T\otimes L$ 截面也有相同系数。三者因此给出同一个理想，而非仅给出相同的几何点集。对任意进一步的态射 $T'\to T$，因子分解经过这个闭子预概形，当且仅当拉回后的这组系数全部为零。

条件 $\alpha=0$ 又恰好等价于可逆商 $E_T\twoheadrightarrow L$ 经由 $Q_T$ 因子分解。由于原商满射，所得 $Q_T\to L$ 仍满射。这正是 $\mathbf P(Q)$ 的函子描述，故给出基变换相容的典范同构
\[
H^{(m)}\simeq\mathbf P_B(Q).
\]
这里所用射影化代表可逆商的约定可参照
\href{https://stacks.math.columbia.edu/tag/01OA}{Stacks Project，27.21节，Tag 01OA}。
本次用上述任意 $T'$ 的因子分解核验了同构，而非只用底域上的点来识别它。该论证没有断言一般拉回后定义理想的生成元成为正则列。

一个可直接检验的非约化例子说明最后的区别。令
\[
R=k[\epsilon]/(\epsilon^2),\qquad P_R=\mathbf P^1_R.
\]
在坐标 $z=Z/W$ 中，取截面 $z=0$，以及余维一线性子预概形 $Z-\epsilon W=0$。它们在每个几何点上相合，但后一个方程沿截面拉回为 $-\epsilon\ne0$。所以这对数据不经过关联预概形；关联预概形在 $\Spec R$ 中的逆像为 $V(\epsilon)$。这一例子及上述理想层论证相符，并未用有限点检验代替一般证明。

\subsection*{四、$X^{(m)}$ 在 $X$ 上的 Grassmann 结构}

设 $f:X\to P$ 任意，并记
\[
L_X=f^*\mathcal O_P(1),\qquad
F_X=\ker(E_X\twoheadrightarrow L_X).
\]
因为右端可逆，这个短正合列局部分裂。故 $F_X$ 局部自由，秩为 $r-1$，且拉回至任意 $T\to X$ 后仍是相应商映射的核 $F_T$。

给定 $X^{(m)}$ 的一个 $T$-值点，等价于给定秩 $m$ 泛商所对应的局部直和因子 $N_T\subset E_T$，并要求 $N_T\subset F_T$。此时有正合列
\[
0\longrightarrow F_T/N_T\longrightarrow E_T/N_T
\longrightarrow L_T\longrightarrow0.
\]
其中中项局部自由，右端可逆，因而此列局部分裂，$F_T/N_T$ 也局部自由。于是 $N_T$ 是 $F_T$ 的秩 $m$ 局部直和因子，等价于 $F_T^\vee$ 的秩 $m$ 局部自由商 $N_T^\vee$。

反过来，给定 $F_T^\vee$ 的秩 $m$ 局部自由商，其对偶给出 $F_T$ 中的局部直和因子 $N_T$。由 $E_T/F_T=L_T$ 以及两次局部分裂，$E_T/N_T$ 局部自由；因此 $N_T$ 也是 $E_T$ 的局部直和因子，给出一个 $B$-值点。它包含在 $F_T$ 中，关联条件自动满足。两个构造互逆，且都与任意 $T'\to T$ 的拉回相容；于是
\[
X^{(m)}\simeq\operatorname{Grass}_m(F_X^\vee)
\]
是 $X$ 上的同构，因而也是法文手添所说的 $P$-同构。这补足旧D17-M009仅以较短文字叙述的反向构造及基变换核验。

原稿的 Euler 式也与此相符。采用可逆商约定时，射影丛上有局部分裂的正合列
\[
0\longrightarrow\Omega^1_{P/S}(1)\longrightarrow E_P
\longrightarrow\mathcal O_P(1)\longrightarrow0.
\]
可在 $E$ 自由的标准图上直接核对：若 $e_0$ 的像生成可逆商，令 $t_i$ 为 $e_i/e_0$ 的坐标，则核由 $e_i-t_i e_0$ 生成；$dt_i$ 与可逆商的生成元相张量对应到这个核基。换图时按 $d(t_i/t_j)$ 计算，正好与该核基的变换相符。这给出所写的全局同构。右端可逆使拉回仍正合，所以
\[
F_X\simeq f^*\bigl(\Omega^1_{P/S}(1)\bigr).
\]
这里不要求 $f$ 平坦；不能因为一般拉回未必左正合，便把这个已经局部分裂的序列误判为需要额外平坦假设。

\subsection*{五、维数、几何整性与两个投影的范围}

令 $N=r-1$。在 $F_X$ 平凡化的开集上，$\operatorname{Grass}_m(F_X^\vee)$ 的标准开图可将商矩阵写成
\[
\begin{bmatrix}I_m&A\end{bmatrix},
\]
其中 $A$ 是任意 $m\times(N-m)$ 矩阵。这给出相对仿射空间开覆盖，各图的相对维数为
\[
m(N-m)=m(r-1-m).
\]
有关商函子及矩阵图的构造可参照
\href{https://stacks.math.columbia.edu/tag/089R}{Stacks Project，27.22节，Tag 089R}。
标准开图从局部生成元中选取一个可逆极大子式得到，因而覆盖整个预概形。由这些开图得到在任意基预概形上的光滑性；它们的数目与维数只使用有限秩条件。

对每个几何纤维，设其域为 $K$。在 $0<m<N$ 时，上述开图均约化；另一方面，$\mathrm{GL}_{N,K}$ 作用在标准商上，得到一个到 Grassmann 预概形的满射态射。这里的满射性可以在任意点的剩余域上选取商与核的基来验证，而不只验证 $K$-有理点。$\mathrm{GL}_{N,K}$ 是仿射空间中行列式非零的整开集，其连续满像不可约。因此该纤维既约化又不可约，即为整概形。一个标准矩阵开图是非空稠密仿射空间，故纤维是有理簇。$m=0$ 或 $m=N$ 时，商函子唯一，纤维为一点；$m>N$ 时为空。这证明正文的非空范围与几何整性，并非仅以单个低维例子猜测一般维数。

法文第66页手添的 $m(r-m)$ 使用了错误的总秩。最小的直接核验是 $r=2,m=1$：$P$ 是射影直线，通过一个点的超平面就是该点，Grassmann 纤维为一点，维数为0。原式给1，订正式给0。这里 $m=r$ 时的空预概形仍可为光滑态射，但其纤维不能称为几何整；故光滑性与非空几何整性不能不分条件地一起断言。

另一投影 $H^{(m)}\to B$ 来自秩 $r-m$ 的 $Q$，在 $0\leq m\leq r-1$ 时为相对维数 $r-m-1$ 的射影丛；$m=r$ 时为空。$X^{(m)}\to B$ 则是线性截面族，不享有同样的普遍结论。例如取 $X$ 为 $\mathbf P^1_k$ 中一个固定的 $k$-点，$m=1$。这时 $X^{(1)}\to X$ 为同构，而 $X^{(1)}\to B\simeq\mathbf P^1_k$ 为一个点的闭浸入。在该参数点取局部坐标 $t$，局部环中乘以 $t$ 为单射，张量到其商 $k$ 后不再单射，因此这一个投影不平坦。它并不否定普遍关联 $H^{(1)}\to B$ 的射影丛结构。

\subsection*{六、迭代记号的限度与核查结论}

作者没有在这里另行定义一个严格的迭代函子。若按先选 $P$ 中的线性子簇、再在其中截取的自然解释来比较，迭代参数保留中间子簇。取 $P=\mathbf P^2$，先取直线再取其上一点，参数为点线旗标；一次取余维2的子簇只参数化点。固定最后的点，中间直线仍有 $\mathbf P^1$ 的选择。因此遗忘映射具有正维纤维，不能被作者试拟的等式视作同构。本版只据此限定符号的解释，没有编造一个原稿未给出的普遍迭代定理或无条件基变换公式。

本模块已经核验的数学内容是：上述有限秩背景下的关联零点结构、两个函子同构、秩和相对维数、光滑及几何整纤维的范围，以及区别两种投影的反例。作者在第66页随后开始对2.1及其他既有命题作推广指示，那部分另有各自的假设和证明依赖；本模块的核验不自动证明它们，也不消除本书其他章节的未决问题。

\subsection*{七、旧D17-M001至D17-M012的逐字存档}

以下逐条保存第十九版旧脚注的原始内容。它们作为历史比较材料保留；其中笼统说“底本”的地方在原制作时主要指英文编辑稿，不能越过上表而当作本次对法文的判断。数学上仍有效的独立说明已在本版正文与本报告中保留或展开；涉及来源归责的改动以D20-12O系列注为准。

\paragraph*{旧注 D17-M001（逐字存档）}
\begin{quote}
这里沿用第11节的 $P^\vee=\mathbf P(E^\vee)$ 及有限秩对偶约定；上标表示余维，下标表示相对维数或商模的秩，须依参数是 $P$ 还是 $E$ 区分。对秩为 $r$ 的 $E$，$0\leq n\leq r-1$ 对应非空线性子簇；$n=r$ 可按上述空子簇约定扩充。
\end{quote}

\paragraph*{旧注 D17-M002（逐字存档）}
\begin{quote}
底本此句写“given dimension $m$”，但前后的上标约定、$\operatorname{Grass}_m(E^\vee)$ 及下一段均使用余维 $m$；故正文更正为“相对余维”。底本也在此从前句秩 $r+1$ 改取秩 $r$，本译文保留这一重新约定，不把两处 $r$ 混为同一秩。
\end{quote}

\paragraph*{旧注 D17-M003（逐字存档）}
\begin{quote}
底本写 $G$ 是 $E_{\operatorname{Gr}}$ 的商，随即说“let us call it $F$”，以后却继续写 $G$。由 $\operatorname{Gr}^m=\operatorname{Grass}_m(E^\vee)$，泛商必来自 $E^\vee$ 而非 $E$；其对偶才给出 $G^\vee\hookrightarrow E$。故补正对偶符号，并将本段泛商统一记为 $G$。底本的临时命名 $F$ 保留在此注，不让它与下一批用于核模的 $F$ 相混。
\end{quote}

\paragraph*{旧注 D17-M004（逐字存档）}
\begin{quote}
底本乘积写成“$P_X\,X\operatorname{Gr}^m$”，且在这一函子说明中混用 $S$ 与 $S'$。这里尚未引入 $X$，二因子均在 $S$ 上，故统一为 $P\times_S\operatorname{Gr}^m$；取值点等价于 $P_{S'}\to S'$ 的截面。
\end{quote}

\paragraph*{旧注 D17-M005（逐字存档）}
\begin{quote}
底本文字把转置后的商写作 $G_T^\vee$，但紧随其后的显示箭头为 $\mathcal O_T(-1)\to G_T$，且下一页截面属于 $G_T(1)$。由有限局部自由模的对偶关系可直接核实正确目标是 $G_T$，正文据此更正。$\mathcal O_T(1)$ 指沿 $T\to P$ 拉回的可逆模；“局部直和因子”来自泛商列的局部分裂。
\end{quote}

\paragraph*{旧注 D17-M006（逐字存档）}
\begin{quote}
这里是概形意义的零点，而不只是点集。局部平凡化 $G_T$ 与 $\mathcal O_T(1)$ 后，三个描述均由同一组 $m$ 个系数生成的理想定义，故具有相同的概形结构。这一说明没有额外断言该组系数在任意基变换后仍构成正则列。
\end{quote}

\paragraph*{旧注 D17-M007（逐字存档）}
\begin{quote}
底本首次引入此记号时漏写下标 $\xi$，后一句明确使用 $Y^m_\xi$，故予统一。显示的纤维积公式只是将该记号的定义展开：若相应线性子簇为 $L_\xi\subseteq P_{S'}$，则 $Y^m_\xi=X_{S'}\times_{P_{S'}}L_\xi$。当 $X\to P$ 不是浸入时，不把它误称为 $P$ 内两个子概形的字面交集。
\end{quote}

\paragraph*{旧注 D17-M008（逐字存档）}
\begin{quote}
底本在此试拟“we might approximately have $X^{(m)(m_1)}=X^{m+m_1}$”。这是关于记号的愿望，不能直接当作两个普遍族的典范同构：连续选取两个线性截面通常保留先后选择或旗标的数据，而一次选取总余维截面只保留最终子空间；二者参数空间一般不同。例如在射影平面中先选直线再选其上一点，参数为点线旗标；直接选余维2子空间只参数化点，忘掉直线的映射具有正维纤维。因此正文只陈述记号意图，原等式保留于此，不把它升格为定理。
\end{quote}

\paragraph*{旧注 D17-M009（逐字存档）}
\begin{quote}
底本写作“$F_X\simeq F^*(\Omega^1_{P/S})(1)$”，把态射与核模的字母混用。这里明确以 $f$ 表示态射，以 $F_X$ 表示核模。射影丛的Euler正合列拉回后仍正合，因为右端可逆且正合列局部分裂。因此 $F_X$ 局部自由，秩为 $r-1$。上述同构也可直接按函子核验：关联条件是泛子模 $G^\vee\subset E_X$ 的像落入 $F_X$；它在 $F_X$ 中仍是局部直和因子，因为 $F_X/G^\vee$ 是满射 $E_X/G^\vee\to\mathcal O_X(1)$ 的局部自由核。对偶后即得到 $F_X^\vee$ 的秩 $m$ 局部自由商，反向构造相同，并与基变换相容。这里的说明是译者补充，不是底本另列的证明。
\end{quote}

\paragraph*{旧注 D17-M010（逐字存档）}
\begin{quote}
底本原式为 $m(r-m)$。但本节上一页已重新规定 $\operatorname{rank}E=r$，所以核模 $F_X$ 的秩是 $r-1$；其秩 $m$ 商的格拉斯曼簇维数为 $m((r-1)-m)$。例如 $r=2,m=1$ 时，射影直线上包含给定点的超平面只有该点，纤维维数为0，不是原式给出的1。若取 $m=r$，该格拉斯曼概形为空，不能无条件声称纤维几何整；故正文明确范围。端点 $m=0,r-1$ 的纤维为一点。中间范围内，矩阵坐标图给出光滑性及有理开图，几何整性是通常格拉斯曼簇的性质。
\end{quote}

\paragraph*{旧注 D17-M011（逐字存档）}
\begin{quote}
这里 $G$ 是前文中的 $E^\vee_{\operatorname{Gr}^{m}}$ 的泛秩 $m$ 商，所以 $G^\vee$ 是 $E_{\operatorname{Gr}^{m}}$ 的局部直和因子，余商的秩为 $r-m$，射影丛相对维数为 $r-m-1$。秩为零的端点 $m=r$ 给出空预概形，故不将其称为非空射影纤维。与前段不同，这里讨论的是 $H^{(m)}\to\operatorname{Gr}^{m}$，而不是 $X^{(m)}\to X$。
\end{quote}

\paragraph*{旧注 D17-M012（逐字存档）}
\begin{quote}
底本在括号中明确保留这一限制。$X^{(m)}=X\times_P H^{(m)}$ 的参数纤维是 $X$ 的线性截面，可能奇异、可约或为空；本译稿不把前段关于另一投影的光滑性误用于这里。
\end{quote}

